\title{xLSTM: Extended Long Short-Term Memory}

\begin{abstract}
During the 1990s, the development of the Long Short-Term Memory (LSTM) architecture was fundamentally shaped by the concepts of gating mechanisms and the constant error carousel. Over time, LSTMs have proven remarkably durable and have underpinned many achievements in deep learning, notably serving as the core of the initial generations of Large Language Models (LLMs). However, the emergence of Transformer architectures—featuring inherently parallelizable self-attention mechanisms—initiated a paradigm shift, ultimately surpassing LSTM-based approaches at large scales. In this work, we revisit LSTMs, asking: To what extent can large-scale language modeling be accomplished by scaling up LSTMs to billions of parameters, incorporating advances from contemporary LLMs, while systematically addressing well-documented shortcomings? Our contributions are two-fold. First, we propose exponential gating, supported by robust normalization and stabilization methods. Second, we redesign LSTM memory components, resulting in (i) the sLSTM, characterized by scalar memory, scalar updates, and revised memory mixing operations, and (ii) the mLSTM, a variant offering complete parallelizability via matrix-based memory and a covariance-driven update strategy. By integrating these modified LSTM elements within residual block frameworks, we construct xLSTM blocks and, by extension, scalable xLSTM architectures through residual stacking. The combination of exponential gating and adapted memory mechanisms enables xLSTM to achieve competitive—or even superior—performance and scaling characteristics when benchmarked against leading Transformer models and State Space Models.\\
Code available at: \url{https://github.com/NX-AI/xlstm}
\end{abstract}

\section{Introduction}
\label{sec:intro}

The Long Short-Term Memory (LSTM) ideas~\citep{Hochreiter:91,Hochreiter:97nips,Hochreiter:97}, i.e., the constant error carousel and gating, were introduced to overcome the vanishing gradient problem of recurrent neural networks~\citep{Hochreiter:91,Hochreiter:00book}:

\begin{align}
\label{eq:lstm_idea}
  c_t \ = \ f_t \ c_{t-1} \ + \ 
          i_t \  z_t \ , \quad  
          h_t \ = \ o_t \ \psi({c_t}) \ . 
\end{align}

The cell state $c_{t-1}$ (shown in green) is incrementally updated via addition with the cell input $z_t$, a process governed by the constant error carousel, and this update is regulated by the input gate $i_t$ and forget gate $f_t$, both mediated through sigmoid activations (depicted in blue). Meanwhile, the output of the memory cell—which corresponds to the hidden state $h_t$—is determined by the output gate $o_t$. Prior to being gated for output, the cell state undergoes normalization or squashing through $\psi$, after which it is transformed into the hidden state.

Long Short-Term Memory networks (LSTMs) have been leveraged across a wide range of disciplines \citep{Hochreiter:01,Hochreiter:07,Schmidhuber:15}, dominating tasks such as text generation until the introduction of Transformer models in 2017~\citep{Vaswani:17}. Their superior performance has been validated in diverse sequence modeling challenges, including but not limited to: text generation~\citep{Graves:13,Karpathy:15blog}, handwriting synthesis~\citep{Graves:13}, machine translation tasks involving sequence-to-sequence models~\citep{Sutskever:14nips}, the assessment of computer programs~\citep{Zaremba:14arxiv}, image captioning~\citep{Karpathy:15,Hossain:19}, program code synthesis~\citep{Karpathy:15blog}, hydrological simulations such as rainfall-runoff modeling~\citep{Kratzert:18,Kratzert:19}, and flood early warning systems~\citep{Nearing:24}. In the domain of reinforcement learning, LSTMs represent the leading sequence modeling architectures and are featured in top-performing systems like AlphaStar for StarCraft II~\citep{Vinyals:17short}, OpenAI Five for Dota 2~\citep{OpenAI:19}, as well as in advanced controllers for nuclear fusion apparatus~\citep{Degrave:22short}. LSTMs demonstrate exceptional capability in capturing and encoding semantic abstractions within their memory components~\citep{Karpathy:15blog}; this capacity is illustrated by the emergence of distinct neurons specializing in numbers, syntax~\citep{Lakretz:19}, linguistic features~\citep{Bau:19}, and sentiment~\citep{Radford:17arxiv}. Despite the advent of newer models, LSTMs continue to find application in critical domains~\citep{Degrave:22short,Nearing:24} and remain a robust choice for sequence modeling challenges.

Although LSTMs have achieved significant advances, they exhibit three principal drawbacks: (i) Inability to update storage choices. This issue is illustrated using the {\it Nearest Neighbor Search} problem (refer to Appendix~\ref{sec:appNearestNeighborSearch}): Given a reference vector, the task is to sequentially examine a sequence to identify the vector most similar to the reference and report its associated value at the conclusion of the sequence. The left panel of Figure~\ref{fig:lstmProblems} presents the mean squared error for this challenge. LSTMs have difficulty revisiting and modifying previously stored values if a closer match is encountered at a later stage, whereas the proposed xLSTM addresses this with exponential gating mechanisms. (ii) Constrained memory capacity; specifically, all information must be compacted into scalar-valued cell states. This shortcoming is demonstrated by {\it Rare Token Prediction}. In the right panel of Figure~\ref{fig:lstmProblems}, the perplexity of token prediction in Wikitext-103~\citep{Merity:17} is plotted against bins representing varying token frequencies. LSTMs tend to underperform on rare tokens, primarily due to their restricted storage capability. xLSTM overcomes this challenge through the use of matrix-based memory. (iii) Inherent sequential nature caused by memory mixing, namely, the presence of recurrent connections between consecutive hidden states that preclude parallel processing. 

These challenges associated with LSTM architectures have opened the door for the adoption of Transformers~\citep{Vaswani:17} within the field of language modeling. This raises the question: by removing these limitations and scaling up LSTM-based models to the scale of contemporary Large Language Models, what levels of language modeling performance can be achieved?

\section{Extended Long Short-Term Memory}
\label{sec:xLSTM}

In order to address the inherent shortcomings of standard LSTMs, the Extended Long Short-Term Memory (xLSTM) framework implements two significant changes to the formulation presented in Equation~\eqref{eq:lstm_idea}. These enhancements---the use of exponential gating and the introduction of new memory arrangements---result in two novel variants within the LSTM framework: (i) the sLSTM (see Section~\ref{sec:sLSTM}), which incorporates a scalar memory, utilizes scalar updates, and employs memory mixing; and (ii) the mLSTM (see Section~\ref{sec:mLSTM}), characterized by a matrix-based memory and an update rule grounded in covariance (outer product) operations that allow complete parallel execution. Both sLSTM and mLSTM improve upon classical LSTM models through the integration of exponential gating mechanisms. The mLSTM achieves parallelizability by forgoing memory mixing, specifically removing hidden-hidden recurrent links. Extensions of both models to multiple memory units are feasible, wherein sLSTM supports memory mixing between cells. Additionally, sLSTM is capable of having several heads, with memory mixing restricted to the cells within a head and not across different heads. This innovation, involving head division and exponential gating in sLSTM, gives rise to a distinct memory mixing approach. For mLSTM, using multiple heads and multiple cells produces similar effects.

Embedding these enhanced LSTM structures into residual block modules gives rise to xLSTM blocks (refer to Section~\ref{sec:xLSTMblock}). When these xLSTM blocks are layered residually within a network, they form xLSTM architectures (see Section~\ref{sec:xLSTMarchitecure}). The overarching design and its modular elements are illustrated in Figure~\ref{fig:xlstm_sketch}.

\subsection{Review of the Long Short-Term Memory}
\label{sec:orgLSTM}

The initial concept of LSTM~\citep{Hochreiter:91,Hochreiter:97nips,Hochreiter:97} presented a scalar memory cell, serving as the main component for information retention and manipulation, designed specifically to mitigate the vanishing gradient issue~\citep{Hochreiter:91,Hochreiter:00book} by utilizing the so-called constant error carousel, i.e., the way in which the cell state is updated. This memory cell architecture features three distinct gating mechanisms—input, output, and forget gates—with the forget gate later proposed by \citet{Gers:00}. The dynamics of the LSTM memory cell at any given time $t$ are governed by the following update rules:

\begin{align}
c_t \ &= \  f_t \ c_{t-1} \ + \ i_t \ z_t &  & &\text{cell state} \\
h_t  \ &= \ o_t \ \tilde{h}_t \ , & \tilde{h}_t \ &= \ \psi \left( c_t\right)
  &\text{hidden state} \\
z_t \ &= \ \varphi \left( \tilde{z}_t \right) \ , 
  &\tilde{z}_t \ &=  \ \mathbf{w}^\top_{z} \ \mathbf{x}_t \ + \
  r_{z}  h_{t-1} \ + \  b_{z} \ \
  &\text{cell input} \\
i_t \ &= \ \sigma \left( \tilde{i}_t  \right) \ , 
  &\tilde{i}_t \ &= \ \mathbf{w}^\top_{i} \ \mathbf{x}_t \ + \
  r_{i}  \ h_{t-1} \ + \  b_{i} \ \
  &\text{input gate} \\
f_t \ &= \ \sigma \left(  \tilde{f}_t \right) \ , 
  &\tilde{f}_t \ &= \ \mathbf{w}^\top_{f} \ \mathbf{x}_t  \ + \
  r_{f}  \ h_{t-1} \ + \  b_{f} \ \
  &\text{forget gate} \\
o_t \ &= \ \sigma \left( \tilde{o}_t \right) \ , 
  &\tilde{o}_t  \ &= \ \mathbf{w}^\top_{o} \ \mathbf{x}_t \ + \
  r_{o}  \ h_{t-1} \ + \  b_{o} \ \
  &\text{output gate} 
\end{align}

The vectors $\mathbf{w}_{z}$, $\mathbf{w}_{i}$, $\mathbf{w}_{f}$, and $\mathbf{w}_{o}$ represent the input weights connecting $\mathbf{x}_t$ to the cell input, input gate, forget gate, and output gate, respectively. Analogously, $r_{z}$, $r_{i}$, $r_{f}$, and $r_{o}$ denote the recurrent weight parameters from the previous hidden state $h_{t-1}$ to the respective components: cell input, input gate, forget gate, and output gate. The bias vectors $b_{z}$, $b_{i}$, $b_{f}$, and $b_{o}$ correspond to these pathways. The cell input and hidden state activations are governed by $\varphi$ and $\psi$ (often set to $\tanh$ functions), with $\psi$ ensuring that the cell state remains bounded. Sigmoid activation functions are applied to all gate components, namely, $\sigma \left( x \right)= 1/(1+\exp(-x))$. Subsequent developments represented several scalar memory cells as a vector $\mathbf{c}_t \in \mathbb{R}^{d}$, facilitating the application of recurrent matrices $\mathbf{R} \in \mathbb{R}^{d \times d}$ to jointly process the outputs of memory cells~\citep{Greff:15}; see Appendix~\ref{sec:apporgLSTM} for further explanation. Ablation experiments revealed that each part of the memory cell is essential for proper functioning~\citep{Greff:15}.

\subsection{sLSTM}
\label{sec:sLSTM}

To enable LSTMs to reconsider and update their storage choices, we incorporate exponential gates (shown in red), alongside mechanisms for normalization and stabilization. Specifically, it is possible for both the input and forget gates to utilize exponential activation functions. For normalization, we propose a dedicated normalizer state, which accumulates the sum resulting from multiplying the input gate by all forthcoming forget gates.

The scalar sLSTM forward pass is:

\begin{align}
\label{eq:slstmforward}
c_t \ &= \  f_t \ c_{t-1} \ + \ i_t \ z_t & & 
 &\text{cell state} \\
n_t \ &= \  f_t \ n_{t-1} \ + \ i_t  & & 
 &\text{normalizer state} \\
h_t \ &= \ o_t \ \tilde{h}_t \ ,
  &\tilde{h}_t \ &= \ c_t / n_t
&\text{hidden state} \\
z_t \ &= \ \varphi \left( \tilde{z}_t \right) \ , 
  &\tilde{z}_t \ &=  \ \mathbf{w}^\top_{z} \ \mathbf{x}_t \ + \
  r_{z}  h_{t-1} \ + \  b_{z}
  &\text{cell input} \\
i_t \ &= \ \!\exp \left( \tilde{i}_t  \right) \ , 
  &\tilde{i}_t \ &= \ \mathbf{w}^\top_{i} \ \mathbf{x}_t \ + \
  r_{i}  \ h_{t-1} \ + \  b_{i}
  &\text{input gate} \\
f_t \ &= \ \sigma \left(  \tilde{f}_t \right) \ \text{OR} \
\!\exp \left(  \tilde{f}_t \right) \ ,
  &\tilde{f}_t \ &= \ \mathbf{w}^\top_{f} \ \mathbf{x}_t  \ + \
  r_{f}  \ h_{t-1} \ + \  b_{f}
  &\text{forget gate} \\
o_t \ &= \ \sigma \left( \tilde{o}_t \right) \ , 
  &\tilde{o}_t  \ &= \ \mathbf{w}^\top_{o} \ \mathbf{x}_t \ + \
  r_{o}  \ h_{t-1} \ + \  b_{o}
  &\text{output gate} 
\end{align}

We transfer the original LSTM gating techniques, i.e., input- and/or hidden-dependent gating plus bias term, to the new architectures. Exponential activation functions can lead to large values that cause overflows. Therefore, we stabilize gates with an additional state \mbox{$m_t$~\citep{Milakov:18arxiv}}:

\begin{align}
\label{eq:slstmstabil}
m_t \ &= \ \max \left( \log ( f_t ) + m_{t-1} , \log ( i_t ) \right) &\text{stabilizer state} \\
i'_t \ &= \ \!\exp \left( \log \left ( i_t \right) - m_t \right) \ = \ \!\exp \left( \tilde{i}_t - m_t \right) \ \, 
  &\text{stabil. input gate} \\
f'_t \ &= \ \!\exp \left( \log \left( f_t \right) + m_{t-1} - m_t \right) \ \, 
  &\text{stabil. forget gate}
\end{align}

As demonstrated in Appendix~\ref{sec:appsLSTM}, substituting $f_t$ with $f'_t$ and $i_t$ with $i'_t$ during the forward computation leaves both the network’s output and the gradients of the loss function with respect to the model parameters unaltered.

\paragraph{New Memory Mixing.} Similar to the conventional LSTM architecture, sLSTM supports the use of several memory cells (refer to Appendix~\ref{sec:appsLSTM} for details). The presence of multiple memory cells allows for memory mixing, implemented through recurrent pathways: $\mathbf{R}_{\mathbf{z}}$, $\mathbf{R}_{\mathbf{i}}$, $\mathbf{R}_{\mathbf{f}}$, $\mathbf{R}_{\mathbf{o}}$, which connect the hidden state vector $\mathbf{h}$ with the memory cell input $\mathbf{z}$ and the gate units $\mathbf{i}$, $\mathbf{f}$, $\mathbf{o}$, respectively. In this approach, a distinctive characteristic of memory mixing arises from the application of exponential gating. The updated sLSTM architecture may feature multiple heads, facilitating memory mixing within a single head but not between different heads. The integration of head structures alongside exponential gating in sLSTM introduces a novel mechanism for mixing information stored in memory.

\subsection{mLSTM}
\label{sec:mLSTM}

To improve the storage capabilities of LSTMs, we upgrade the traditional scalar memory cell $c \in \mathbb{R}$ to a higher-dimensional matrix representation $\mathbf{C} \in \mathbb{R}^{d \times d}$. Consequently, information retrieval is accomplished through matrix multiplication. At each time step $t$, the model aims to encode a key–value pair consisting of a key vector $\mathbf{k}_t \in \mathbb{R}^d$ and a value vector $\mathbf{v}_t \in \mathbb{R}^d$ (utilizing nomenclature from Transformer models). Subsequently, at a future time $t + \tau$, a query vector $\mathbf{q}_{t+\tau} \in \mathbb{R}^d$ is used to recover the original value $\mathbf{v}_t$. This scenario aligns with the classical framework of Bidirectional Associative Memories (BAMs) \citep{Kohonen:72,Anderson:72,Nakano:72,Anderson:77}.

The covariance update rule~\citep{Sejnowski:77,Dayan:91} for storing a key-value pair is

\begin{align}
\mathbf{C}_t \ &= \ \mathbf{C}_{t-1} \ + \ \mathbf{v}_{t} \ \mathbf{k}_t^\top \ .
\end{align}

We presume that layer normalization is applied prior to projecting inputs into keys and values, resulting in these representations having zero mean. The covariance update method provides an optimal mechanism~\citep{Dayan:91} for achieving maximum separability among retrieved binary vectors, a condition that corresponds to optimizing the signal-to-noise ratio. Greater separability can be realized by focusing retrieval on pairwise interactions, which accepts the quadratic computational cost similar to that of attention mechanisms~\citep{Krotov:16,Krotov:17,Ramsauer:21}. The covariance update approach is also functionally equivalent to Fast Weight Programmers~\citep{Schmidhuber:92ncfastweights,Schlag:21}, which have subsequently incorporated a constant decay factor for $\mathbf{C}_{t-1}$ and a fixed learning rate for scaling 
$\mathbf{v}_{t} \mathbf{k}_t ^\top$~\citep{Ba:16}.
Following this analogy, we adapt the covariance update method within the LSTM architecture, mapping the forget gate to the decay factor, the input gate to the learning rate, and using the output gate to modulate the retrieved output.

For this matrix memory, the normalizer state is the weighted sum of key vectors, where each key vector is weighted by the input gate and all future forget gates. Again, the normalizer state keeps record of the strength of the gates. Since the dot product between query and normalizer state can be close to zero, we use the absolute value of this dot product and lower bound it by a threshold (typically 1.0) as done previously~\citep{Sun:23arxiv}. The mLSTM forward pass is:

\begin{align}
\label{eq:mlstm_recurrent_begin}
\mathbf{C}_t \ &= \  f_t \ \mathbf{C}_{t-1} \ + \ 
  i_t \ \mathbf{v}_t \ \mathbf{k}_t^\top & &  &\text{cell state} \\
\mathbf{n}_t \ &= \  f_t \ \mathbf{n}_{t-1} \ + \ 
  i_t \ \mathbf{k}_t & &  &\text{normalizer state} \\
\mathbf{h}_t  \ &= \ \mathbf{o}_t \ \odot \ \tilde{\mathbf{h}}_t
\ , \qquad \qquad 
\tilde{\mathbf{h}}_t \ = \   \mathbf{C}_t \mathbf{q}_t \ / \ 
  \max \left\{ |\mathbf{n}_t^\top \mathbf{q}_t|, 1 \right\} 
   &&&\text{hidden state} \\
\mathbf{q}_t \ &= \ \mathbf{W}_q \ \mathbf{x}_t \ + \ \mathbf{b}_q  & &  &\text{query input} \\
\mathbf{k}_t \ &= \ \frac{1}{\sqrt{d}} \mathbf{W}_k \ \mathbf{x}_t \ + \ \mathbf{b}_k  & &  &\text{key input} \\
\mathbf{v}_t \ &= \ \mathbf{W}_v \ \mathbf{x}_t \ + \ \mathbf{b}_v  & &  &\text{value input} \\
i_t \ &= \ \!\exp \!  \! \left( \tilde{i}_t  \right) \ , 
 \qquad \qquad \ \ \ \ \,
  \tilde{i}_t \ = \ \mathbf{w}^\top_{i} \ \mathbf{x}_t \ + \  b_{i} \ \ \ 
  &&&\text{input gate} \\
f_t \ &= \ \sigma \!  \left(  \tilde{f}_t \right) \ \text{OR} \ \!\exp \!  \! \left(  \tilde{f}_t \right) , \ \ \: \!
\tilde{f}_t \ = \ \mathbf{w}^\top_{f} \ \mathbf{x}_t  \ + \
  b_{f}
  &&& \text{forget gate} \\
\label{eq:mlstm_recurrent_end}
\mathbf{o}_t \ &= \ \sigma \left( \tilde{\mathbf{o}}_t \right) \ , \qquad \qquad \quad \ \ \ \,
  \tilde{\mathbf{o}}_t  \ = \ \mathbf{W}_{\mathbf{o}} \ \mathbf{x}_t \ + \
  \mathbf{b}_{\mathbf{o}}
  &&&\text{output gate} 
\end{align}

mLSTM, similar to standard LSTM, is capable of incorporating several memory cells. In mLSTM, having multiple cells or multiple heads results in the same functionality due to the absence of interactions between memories. To ensure stable operation of the exponential gates within mLSTM, we apply the stabilization methods employed for sLSTM (refer to Equation~\ref{eq:slstmstabil}). Because there is no mixing across memories in the mLSTM, its recurrence relations can be recast in a parallel form. Further elaboration can be found in Appendix~\ref{sec:appmLSTM}.

\subsection{xLSTM Architecture}
\label{sec:xLSTMarch}

\paragraph{xLSTM Blocks.}
\label{sec:xLSTMblock}
The purpose of an xLSTM block is to encode the historical sequence information through non-linear transformations within a high-dimensional representation, thus enhancing the ability to distinguish between distinct historical contexts. Effectively isolating past contexts is fundamental for the accurate prediction of forthcoming sequence elements, such as the subsequent token. Our approach is grounded in Cover’s Theorem~\citep{Cover:65}, which asserts that patterns, when projected into a higher-dimensional space via non-linear embeddings, become more amenable to linear separation compared to their original form. To this end, we investigate two kinds of residual block designs: (i) A residual block featuring a post up-projection (analogous to Transformers), wherein the past is first non-linearly processed in the original feature space. Afterwards, a linear projection increases the dimensionality, followed by a non-linear activation, and a final linear projection brings the representation back to the initial dimension; see the left part of Figure~\ref{fig:Backbones}
and the third column of Figure~\ref{fig:xlstm_sketch}. 
A comprehensive illustration is provided in Figure~\ref{fig:appsLSTM_detailed}
in the appendix. (ii) A residual block utilizing pre up-projection (reminiscent of State Space Models), where the mapping to a higher-dimensional space is performed linearly at the outset,

The approach non-linearly condenses historical information within a high-dimensional representation before applying a linear transformation to restore it to the original space. In the case of an xLSTM block that incorporates an sLSTM module, we employ the post up-projection block. Conversely, for an xLSTM block built around an mLSTM, the pre up-projection block is chosen, owing to the increase in memory capacity enabled by the higher-dimensional embedding. For further explanation, please consult the left panel in Figure~\ref{fig:Backbones}, the third column of Figure~\ref{fig:xlstm_sketch}, or refer to Figure~\ref{fig:appsLSTM_detailed} in the appendix.

\paragraph{xLSTM Architecture. }
\label{sec:xLSTMarchitecure}
The xLSTM architecture is devised by stacking fundamental units in a residual configuration, as described in prior works~\citep{Srivastava:15,He:16}. Our implementation builds upon the widely adopted pre-LayerNorm~\citep{Ba:16b} residual backbone structure prevalent in modern Large Language Models. This is illustrated in the last column of Figure~\ref{fig:xlstm_sketch}.

\subsection{Memory and Speed Considerations}

In contrast to Transformers, xLSTM architectures exhibit linear computational complexity and maintain constant memory usage relative to sequence length. Thanks to its compressive memory design, xLSTM is particularly advantageous for industrial use cases and deployment on edge devices.

The mLSTM's memory is parameter-free but involves a computational burden, as it relies on $d \times d$ matrix storage and updates. This design represents a compromise between available memory capacity and computational cost. However, since these operations are amenable to parallelization on GPUs, their impact on actual execution time is minimal.

Similar to FlashAttention~\citep{Dao:22,Dao:23} and GLA~\citep{Yang:23arxiv}, mLSTM computations are parallelizable. By contrast, sLSTM cannot be parallelized efficiently due to memory mixing from hidden-to-hidden interactions. Nevertheless, our optimized CUDA implementation—tuned at the GPU register level—achieves sLSTM execution speeds generally within a factor of two of those attained by mLSTM.

\section{Related Work}
\label{sec:related}

\paragraph{Linear Attention.} A variety of approaches have been proposed to reduce the Transformer's quadratic computational complexity with respect to context length, aiming to achieve linear scalability. For instance, the Synthesizer generates artificial attention weights independently of direct token-to-token relationships~\citep{Tay:20arxiv}. Linformer employs a low-rank factorization to express self-attention, enabling a linear approximation in computation~\citep{Wang:20arxiv}. The Linear Transformer restructures the attention procedure to obtain a linear form~\citep{Katharopoulos:20}. Performer approximates the attention's softmax operation using positive orthogonal random features, thus achieving linear complexity~\citep{Choromanski:21}. Additionally, mechanisms such as Structured Global Convolution (SGConv)~\citep{Li:22} and the Hyena Hierarchy~\citep{Poli:23} substitute attention with rapid, long-range convolutional operators.

\paragraph{State Space Models.} State Space Models (SSMs) have gained significant attention recently, largely due to their linear scalability with respect to sequence length and their competitive results in comparison to Transformers. Among the earliest in this line of work was the Structured State Space sequence model (S4)~\citep{Gu:21}. This was soon succeeded by models such as the Diagonal State Space (DSS) model~\citep{Gupta:22}, Gated State Space (GSS) models~\citep{Mehta:22}, S5 model~\citep{Smith:22}, Bidirectional Gated SSM (BiGS)~\citep{Wang:22}, H3 model~\citep{Fu:23}, and more recently, Mamba~\citep{Gu:24arxiv}.

\paragraph{Recurrent Neural Networks.} In response to the challenge of handling longer context lengths efficiently, Recurrent Neural Networks (RNNs) have emerged as viable alternatives to the Transformer architecture and its associated attention mechanism, owing to their linear computational complexity with respect to context. Recent advancements such as Deep Linear Recurrent Units (LRUs) have yielded strong outcomes in language modeling tasks~\citep{Orvieto:23, De:24arxiv}. Similarly, models like Hierarchically Gated Linear RNN (HGRN)~\citep{Qin:23} and its successor HGRN2~\citep{Qin:24arxiv} have also demonstrated promising capabilities. Among notable RNN-based approaches to large language modeling, RWKV~\citep{Peng:23arxivshort,Peng:24arxivshort} stands out, offering performance levels that rival those of Transformer-based models.

\paragraph{Gating.}
The concept of gating, originally introduced by LSTM, has become central in a variety of more recent architectures, often re-examined or reframed. Gating mechanisms have been integrated into HGRN~\citep{Qin:23}, HGRN2~\citep{Qin:24arxiv}, Gated Linear Attention (GLA)~\citep{Yang:23arxiv}, Gated State Space (GSS) models~\citep{Mehta:22}, Bidirectional Gated SSM (BiGS)~\citep{Wang:22}, Moving Average Equipped Gated Attention (MEGA)~\citep{Ma:22}, as well as models such as RWKV~\citep{Peng:23arxivshort} and Mamba~\citep{Gu:24arxiv}.

\paragraph{Covariance Update Rule.} To improve the memory capabilities of our model, we augmented the mLSTM cell with a matrix-based memory incorporating a covariance update rule. This type of update approach is also used by other models, including Fast Weight Programmers~\citep{Schmidhuber:92ncfastweights,Schlag:21}, RWKV-5 and RWKV-6~\citep{Peng:24arxivshort}, Retention~\citep{Sun:23arxiv}, Linear Transformer~\citep{Katharopoulos:20}, and HGRN2~\citep{Qin:24arxiv}.

\paragraph{Most Related.} The models most akin to xLSTM in terms of underlying principles include Retention~\citep{Sun:23arxiv}, RWKV~\citep{Peng:23arxivshort,Peng:24arxivshort}, and HGRN2~\citep{Qin:24arxiv}, each of which implements matrix memory and/or employs gating mechanisms. Notably, these methods differ from the recently introduced sLSTM as they do not facilitate memory mixing. The capability for memory mixing is crucial for addressing state tracking tasks, thereby rendering LSTMs more expressive than both State Space Models (SSMs) and Transformers~\citep{Merrill:24,Deletang:23}. State tracking is particularly important for tasks such as code evaluation or maintaining continuity of entities across extensive narratives.

\paragraph{Residually Stacking Architectures.} In alignment with the design of virtually all state-of-the-art large deep neural networks, xLSTM models are assembled through the residual stacking of modular building blocks~\citep{Srivastava:15,He:16}.

Such architectural strategies have paved the way for the advancement of deep convolutional networks~\citep{He:16} as well as Transformers~\citep{Vaswani:17}. Transformers serve as the foundational mechanism powering major Large Language Models (LLMs), including GPT-3~\citep{Brown:20short}, ChatGPT~\citep{Schulman:22short}, GPT-4~\citep{Achiam:23gpt4short}, Megatron-LM~\citep{Shoeybi:19}, Gopher~\citep{Rae:21short}, ERNIE 3.0 Titan~\citep{Wang:21arxivshort}, GLaM~\citep{Du:21glamshort}, Chinese M6~\citep{Lin:21}, multilingual AlexaTM 20B~\citep{Soltan:22}, OPT~\citep{Zhang:22opt}, Chinchilla~\citep{Hoffmann:22short}, BLOOM~\citep{Scao:22arxivshort}, GLM-130B~\citep{Zeng:22arxivshort}, LaMDA~\citep{Thoppilan:22short}, PaLM~\citep{Chowdhery:22short}, Llama~\citep{Touvron:23arxiv}, and Gemini~\citep{GeminiTeam:23arxiv,Reid:24arxiv}.

\section{Experiments}
\label{sec:Exp}

In this section, we present an empirical study of xLSTM, positioning it relative to prevailing approaches, with emphasis on the task of language modeling. Section~\ref{sec:ExpSyntethic} is devoted to analyzing the unique abilities of xLSTM on a series of synthetic benchmarks. 
Subsequently, in Section~\ref{sec:ExpComparison}, we systematically compare different state-of-the-art language modeling models by measuring validation perplexity, utilizing a training corpus of 15B tokens from SlimPajama~\citep{Soboleva:23}. We supplement these findings with ablation experiments targeting components of xLSTM on the identical dataset. Furthermore, we investigate the models’ scaling properties following approaches similar to those of \citet{Kaplan:20} and \citet{Brown:20short}.
In Section~\ref{sec:ExpLanguage}, 
a more extensive set of language modeling experiments is provided. We contrast xLSTM with the top-performing models identified in Section~\ref{sec:ExpComparison}, after training on a substantially larger corpus of 300B tokens from SlimPajama~\citep{Soboleva:23}. Our analysis proceeds in four stages: first, we evaluate the models’ generalization to longer input sequences; second, we benchmark validation perplexity alongside outcomes on downstream tasks as discussed in~\citep{Sutawika:24}; third, we measure their efficacy on the PALOMA language benchmark, covering 571 distinct textual domains~\citep{Magnusson:23arxivshort}; and fourth, we re-examine scaling behaviors, now under substantially increased training data by a factor of 20.

\label{sec:xlstm_blocks_notation}
Throughout our experiments, the designation xLSTM[$a$:$b$] is employed to indicate a proportion of $a/b$ between mLSTM-based and sLSTM-based xLSTM blocks. As an illustration, xLSTM[7:1] denotes a configuration where, from a total of eight blocks, seven utilize mLSTM, while one employs sLSTM. With a typical overall block count of 48, this allocation corresponds to 42 blocks based on mLSTM and 6 on sLSTM. Additionally, across all experimental settings, we incorporate up-projection blocks prior to mLSTM blocks (pre) and following sLSTM blocks (post), respectively.

\subsection{Synthetic Tasks and Long Range Arena}
\label{sec:ExpSyntethic}
We begin by evaluating the impact of xLSTM’s exponential gating mechanism with memory mixing on formal language tasks~\citep{Deletang:23}. Next, we investigate how well xLSTM’s matrix memory architecture performs on the Multi-Query Associative Recall task~\citep{Arora:23arxiv}. Lastly, we examine xLSTM’s ability to handle lengthy sequences using benchmarks from the Long Range Arena~\citep{Tay:21}.

\paragraph{Evaluation of xLSTM’s Exponential Gating with Memory Mixing.} We assess the capabilities of xLSTM’s exponential gating combined with memory mixing, a mechanism designed to address state tracking challenges~\citep{Merrill:24, Merrill:23}. To facilitate evaluation, we adopt and expand upon the formal language tasks from \citet{Deletang:23}, adapting them to support multi-length training, which is necessary for investigating length extrapolation. Comprehensive task descriptions and additional empirical findings are presented in Appendix~\ref{sec:appExpSynthetic-formalLang}. For comparative analysis, we include other prominent architectures such as Transformers, State Space Models, and traditional Recurrent Neural Networks. The methods are assessed based on accuracy, measured specifically on task-relevant tokens and normalized from 0 (chance) to 1 (ideal performance). Our experiments involve 2-block variants of the architectures: xLSTM[0:1] (sLSTM component only), xLSTM[1:0] (mLSTM component only), xLSTM[1:1], alongside Llama, Mamba, RWKV, Retention, Hyena, LSTM, and LSTM incorporated within Transformer blocks (LSTM (Block)). The comparative results are provided in Figure~\ref{fig:formal-main}. Architectures lacking memory mixing—such as Transformers or State Space Models—are unable to resolve tasks that depend on state tracking, including basic regular grammars exemplified by the parity task. This finding is consistent with prior work showing that RNNs have greater expressive capacity for such tasks compared to Transformers and State Space Models~\citep{Merrill:24, Merrill:23, Deletang:23}.

\paragraph{Test of xLSTM's Memory Capacities on Associative Recall Tasks.} In this study, we examine the novel matrix memory mechanism of xLSTM by assessing its ability to retain information in the Multi-Query Associative Recall task~\citep{Arora:23arxiv}. During this task, sequences are constructed from randomly sampled key-value pairs drawn from an extensive vocabulary; these associations must be stored for subsequent retrieval. To further challenge the models beyond the original benchmark, we scale the number of key-value pairs as high as 256 and lengthen the input context to 2048. This allows for a more comprehensive investigation of the memory limitations of each architecture. We benchmark 2-block model variants, including Llama, Mamba, RWKV-5, RWKV-6, xLSTM[1:1], and xLSTM[1:0], measuring their retrieval accuracy as a performance metric. Transformers such as Llama serve as a point of reference, given their exponentially large memory in relation to the coding dimension~\citep{Ramsauer:21}, making them the benchmark for this task. Figure~\ref{fig:mqar-main} presents the comparative outcomes. Among the non-Transformer candidates, xLSTM[1:1] consistently achieves the highest recall accuracy, even in smaller model configurations. Notably, incorporating an sLSTM block does not reduce memory retention but in fact enhances it, particularly evident when handling the most challenging setting with 256 pairs. Further data, available in Appendix~\ref{sec:appExpSynthetic-mqar}, support these results and show that the superior memory capability of xLSTM also enables successful extrapolation to longer contexts beyond the training regime.

\paragraph{Test of xLSTM's Long Context Capabilities on Long Range Arena.} In order to evaluate how well xLSTM manages extended sequences and substantial contexts, we benchmark various approaches using the Long Range Arena~\citep{Tay:21}. Results indicate that xLSTM consistently achieves high performance across all evaluated tasks, highlighting the architecture's notable effectiveness in addressing diverse challenges inherent to long context scenarios.

For more details, see Appendix~\ref{sec:appExpSynthetic-lra}.

\subsection{Method Comparison and Ablation Study}
\label{sec:ExpComparison}

In order to investigate our central research question—specifically, the capabilities of our novel LSTM variants at larger scales for language modeling—we conduct experiments by training xLSTMs, Transformers, State Space Models, along with various other approaches, on a dataset of 15B tokens sourced from SlimPajama, all within a consistent auto-regressive framework. Following training, we assess model performance on the validation set and carry out ablation analyses focused on the xLSTMs.

\paragraph{Comparison of xLSTM with Alternative Approaches.} To facilitate comparison, we trained a suite of models on a dataset of 15B tokens sourced from SlimPajama~\citep{Soboleva:23}. Model quality was assessed using perplexity measured on the validation partition. The suite comprised xLSTM (the proposed architecture), GPT-3 (Transformer)~\citep{Brown:20short}, Llama (Transformer)~\citep{Touvron:23arxiv}, H3 (SSM)~\citep{Fu:23}, Mamba (SSM)~\citep{Gu:24arxiv}, RWKV-4 (RNN)~\citep{Peng:23arxivshort}, RWKV-5 (RNN)~\citep{Peng:24arxivshort}, RWKV-6 (RNN)~\citep{Peng:24arxivshort}, GLA (linear Transformer)~\citep{Yang:23arxiv}, HGRN (RNN)~\citep{Qin:23}, HGRN2 (RNN)~\citep{Qin:24arxiv}, RetNet (linear Transformer)~\citep{Sun:23arxiv}, Hyena (linear Transformer)~\citep{Poli:23}, along with xLSTM[1:0] and xLSTM[7:1] (cf. Section~\ref{sec:xlstm_blocks_notation}). All models were trained using mixed-precision arithmetic, except RWKV-5, RWKV-6, GLA, and HGRN2, where PyTorch's native automated mixed-precision was not leveraged (see also Appendix Section~\ref{sec:appGenTrainProc}). Methods were grouped as (a) Transformers, (b) State Space Models (SSMs), and (c) Recurrent Neural Networks (RNNs) as well as linear Transformers, with the latter substituting the canonical attention mechanism with linear methods. For parity, all model architectures matched GPT-3’s parameter budget of 350M, with an embedding dimension of 1024 and 24 residual layers. GPT-3 is unique among these models in sharing its token and output embedding matrices, resulting in a reduced parameter count. As indicated by Table~\ref{tab:spaj15b_model_results}, xLSTM achieves the lowest validation perplexity across all evaluated approaches. Further information is provided in Appendix~\ref{sec:appExpComparison}. Figure~\ref{fig:temp_sclaw15B} depicts the scaling trends, illustrating xLSTM’s superior expected performance for even larger model scales.

\paragraph{Ablation Studies.}
Table~\ref{tab:spaj15b_model_results} and Figure~\ref{fig:temp_sclaw15B} illustrate that xLSTM attains strong performance in language modeling tasks after being trained on 15B SlimPajama tokens. To systematically analyze the individual modifications from LSTM to xLSTM, we incrementally transform a standard LSTM model towards the full xLSTM design. The process begins with embedding the LSTM layers within pre-LayerNorm residual backbones. Next, we adapt this framework by implementing a post up-projection block. The final stage incorporates exponential gating as well as the matrix memory component.

The results are shown in Table~\ref{tab:ablstudies} (top). 

Ablation experiments identify both the exponential gating and the matrix memory as key contributors to the observed performance gains. Given the critical role of gating in RNNs and State Space Models, we further investigate various gating strategies. Results presented in Table~\ref{tab:ablstudies} (bottom) demonstrate that making each gate both learnable and responsive to the input leads to progressive improvements. Further analyses of specific backbone elements are detailed in Appendix~\ref{sec:appAblDetails}.

\subsection{xLSTM as Large Language Model}
\label{sec:ExpLanguage}

This work concludes with extensive language modeling experiments designed to assess the viability of xLSTM as a large language model (LLM). To this end, we scale up the training regime, utilizing 300B tokens sourced from the SlimPajama dataset, which aligns with the dataset sizes adopted in studies such as Mamba~\citep{Gu:24arxiv} and Griffin~\citep{De:24arxiv}.

Our experimental analysis pits xLSTM against RWKV-4, Llama, and Mamba, each exemplifying one of the methodological categories introduced in Section~\ref{sec:ExpComparison}. RWKV-4 serves as our RNN benchmark due to the fact that suitable training precision configurations for RWKV-5, RWKV-6, and HGRN2 were not established until after the commencement of the 300B token experiments (see Appendix~\ref{sec:appGenTrainProc}).

We conduct training runs across multiple model capacities (125M, 350M, 760M, 1.3B), subject each to evaluations of their ability to extrapolate to longer sequences, and measure their efficacy on the corresponding validation sets. Further, we analyze downstream task performance, conduct comprehensive language modeling evaluations using the 571 domains provided by the PALOMA benchmark, and, ultimately, explore the models' adherence to established scaling laws.

\paragraph{Sequence Length Extrapolation.}
\label{sec:spaj300B_ctx_extrapolation}
To begin, we examine how well large-scale models—specifically, xLSTM, RWKV-4, Llama, and Mamba at the 1.3B parameter scale—can generalize to sequences longer than those seen during training. Each model is trained using a fixed context length of 2048 tokens and then evaluated on contexts extending up to 16384 tokens. The resulting performance is depicted in Figure~\ref{fig:spaj300B_seq_extrapolation}. Notably, xLSTM displays a consistent ability to sustain low perplexity scores even as the tested context length increases, setting it apart from the other models.

\paragraph{Validation Perplexity and Downstream Tasks.}
\label{sec:spaj_downstream_eval}
Next, we assess the capabilities of xLSTM, RWKV-4, Llama, and Mamba across all parameter sizes by measuring their validation perplexity on the SlimPajama dataset for next-token prediction, as well as their effectiveness on a selection of downstream benchmarks aimed at evaluating common sense reasoning skills.

The third column in Table~\ref{tab:lmeval} presents the perplexity scores on the validation set for various methods. The models xLSTM[1:0] and xLSTM[7:1] consistently achieve the lowest perplexity across all model scales. The subsequent columns in Table~\ref{tab:lmeval} summarize the results on downstream tasks. In nearly every scenario and for each model size, xLSTM outperforms competing methods; exceptions occur primarily for the ARC task, where Mamba occasionally attains the best scores. For further information, refer to Appendix~\ref{sec:appExpLanguage}.

\paragraph{Performance on PALOMA Language Tasks.} Thirdly, we evaluate the capability of xLSTM, RWKV-4, Llama, and Mamba, across all model sizes, on PALOMA language tasks~\citep{Magnusson:23arxivshort} using next token prediction. The assessment utilizes perplexity as the metric, covering 571 text domains ranging from sources like nytimes.com to Reddit communities such as r/depression. 

Table~\ref{tab:ppleval} displays the measured perplexity, organizing the results by language modeling tasks (first seven columns) and a set of fine-grained domain benchmarks (last five columns). For these language tasks, xLSTM[1:0] outperforms xLSTM[7:1]. 

Across these domains, xLSTM[1:0] achieves lower perplexity than Mamba in 568 out of 571 (99.5\%) cases, outperforms Llama in 486 out of 571 (85.1\%) domains, and is superior to RWKV-4 in 570 out of 571 (99.8\%) domains. Further breakdown of these results can be found in Appendix~\ref{sec:appExpLanguage}.

\paragraph{Scaling Laws.} Fourth, we investigate the scaling properties characterized by a power-law relationship, which enable the projection of model performance at increased parameter counts~\citep{Kaplan:20,Brown:20short}. In Figure~\ref{fig:temp_sclaw300B}, we illustrate these scaling dynamics. Although all evaluated models exhibit analogous scaling trends, they are differentiated by varying offsets. Among them, RWKV-4 shows the least favorable results, trailed by Llama and then Mamba. xLSTM surpasses Mamba, exhibiting a comparable improvement over Mamba as Mamba displays over Llama. This observed scaling behavior suggests that as model size increases, xLSTM is likely to maintain superior performance in comparison with both Transformer-based and State-Space models.

\textbf{Generation Times and Maximal Throughput.} In our evaluation, we present the text generation latency in Figure~\ref{fig:inference_time_generation} and report the peak decoding throughput in Figure~\ref{fig:inference_time_decoding_throughput} (left) across our xLSTM variants at the 1.3B parameter scale. For comparison, we include HuggingFace implementations of Mamba, Llama, and RWKV of comparable size, with Llama utilizing a static key-value cache. It is important to note that, at the time these experiments were conducted, neither Transformer Model cache compilation nor the use of \texttt{torch.compile} for Mamba was functioning. All text generation benchmarks were performed using batch size 1 and a pre-fill value of 16 tokens for each model—settings particularly favorable for the Transformer architecture.

From Figure~\ref{fig:inference_time_generation}, it is apparent that xLSTM, alongside the other recurrent architectures such as Mamba and RWKV-4, exhibits linear scaling behavior, in contrast to the quadratic time complexity observed with Llama. To analyze throughput, we assessed Llama's performance over varying batch sizes and pre-fill lengths. The findings depicted in Figure~\ref{fig:inference_time_decoding_throughput} (right) reveal that, by virtue of its constant memory footprint, xLSTM is able to sustain significantly larger batch sizes compared to Llama and consequently delivers the highest throughput.

\section{Limitations}

(i) Unlike mLSTM, the architecture of sLSTM enforces sequential operations for memory mixing, thus limiting opportunities for parallel computation and making high-speed parallel implementations infeasible. Nonetheless, we have designed an efficient CUDA kernel for sLSTM, which achieves execution times currently less than double those of our parallelized mLSTM version. (ii) It is important to note that our mLSTM CUDA kernels have not been subjected to advanced optimization, and as a result, our current implementation operates at roughly one-quarter the speed of frameworks like FlashAttention or the scan routines in Mamba. There is substantial potential to accelerate mLSTM with more sophisticated CUDA kernels akin to those in FlashAttention. (iii) The use of $d \times d$ matrices in mLSTM increases computational requirements. However, since memory update and readout are parameter-free and readily parallelizable with standard matrix computations, the extra computational complexity does not lead to a major increase in total run time. (iv) The choice of initialization for the forget gates is critical and warrants careful consideration. (v) The independence of the matrix memory from the sequence length means that expanding the sequence could surpass the capacity of the memory for particularly long contexts. However, in practice, this has not posed an issue for sequence lengths up to 16k, as discussed in Section~\ref{sec:spaj300B_ctx_extrapolation}. (vi) Given the intensive computational demand of large-scale language model experiments, we were unable to conduct a thorough optimization of either the architecture or the hyperparameters, especially in the case of larger xLSTM models. We believe that comprehensive tuning efforts will be necessary for xLSTM architectures to realize their optimal performance.

\section{Conclusion}
Our investigation has provided a partial response to our initial query: How effective is language modeling when LSTM architectures are scaled to the order of billions of parameters? Our findings so far suggest that the answer is: ``At least as far as prevailing architectures, such as Transformers and State Space Models''. Through the integration of exponential gating, memory mixing, and an innovative memory structure, we have advanced LSTM into the xLSTM variant. Empirical results demonstrate that xLSTM achieves competitive, and often superior, performance in language modeling tasks relative to leading methods including Transformers and State Space Models. The observed scaling laws further imply that further enlarging xLSTM models positions them as strong contenders to the current generation of Large Language Models that predominantly utilize Transformer architectures. Furthermore, xLSTM shows promise for significantly influencing other domains such as Reinforcement Learning, Time Series Prediction, and the simulation of physical systems.

\end{document}